Ce projet de fin d'études présente la conception et la réalisation de \textbf{PharmacieConnect}, une
plateforme web et mobile liée aux pharmacies de garde en Tunisie. L'idée de départ est simple, mais elle
devient vite délicate dans la pratique : le citoyen doit trouver une information fiable au bon moment, tandis
que l'administration doit pouvoir corriger les pharmacies, les plannings de garde et les données associées
sans perdre le contrôle sur leur qualité. Dans notre cas, cette contrainte a orienté la solution vers deux
usages complémentaires : la consultation mobile et la gestion sécurisée des données.

La partie serveur a été développée avec \textbf{FastAPI}. Elle centralise les données, les règles métier, les
endpoints REST ainsi que les mécanismes de sécurité basés sur les jetons JWT, les rôles et les périmètres
régionaux. Autour de ce backend, un tableau de bord \textbf{React}/\textbf{Vite} permet aux utilisateurs
autorisés de suivre et modifier les pharmacies, les gardes, les médicaments, les imports CSV, les indicateurs
et les journaux d'audit. L'application mobile, réalisée avec \textbf{Expo} et \textbf{React Native}, reprend
les besoins côté citoyen : recherche par texte, consultation selon la date, géolocalisation, catalogue des
médicaments et accès à un assistant conversationnel de premiers secours.

La réalisation ne s'est pas limitée aux interfaces. Elle a aussi demandé un travail de modélisation avec
\textbf{SQLAlchemy}, de validation des entrées, de contrôle d'accès par rôle et de journalisation des actions
sensibles. Nous avons séparé les interfaces clientes, les services métier et la couche de persistance afin de
garder une architecture plus lisible pendant les tests backend et les corrections. L'assistant de premiers
secours utilise, quant à lui, une récupération d'informations hybride. Les réponses restent volontairement
prudentes, car certaines situations nécessitent une prise en charge médicale et ne doivent pas être traitées
comme de simples questions applicatives.

\textbf{Mots-clés :} pharmacie de garde, FastAPI, React, Vite, Expo, React Native, JWT, API REST,
SQLAlchemy, tableau de bord, application mobile, assistant de premiers secours.
